\documentclass{book}
\usepackage[english]{babel}
\usepackage{amsmath,amssymb}
\usepackage{graphicx}
\usepackage{geometry}
\geometry{a4paper, margin=1in}
\usepackage{xcolor}
\usepackage{framed}
\usepackage{titlesec}
\usepackage{tabularx}
\usepackage{booktabs}
\usepackage{hyperref}

\setlength{\parindent}{0pt}
\setlength{\parskip}{1ex}

\definecolor{everflame}{RGB}{180, 40, 40}
\definecolor{templelight}{RGB}{240, 200, 100}
\definecolor{chainlantern}{RGB}{100, 100, 120}
\definecolor{marginalia}{RGB}{80, 80, 80}

\newcommand{\citation}[1]{\textcolor{marginalia}{\small\textit{#1}}}
\newcommand{\ritetag}[1]{\textsc{#1}}
\newcommand{\patron}[1]{\textcolor{everflame}{\textsc{#1}}}
\newcommand{\order}[1]{\textcolor{templelight}{\textsc{#1}}}
\newcommand{\danger}[1]{\textcolor{red}{\textbf{#1}}}

\titleformat{\chapter}{\huge\bfseries}{Chapter \thechapter}{20pt}{\huge}
\titleformat{\section}{\Large\bfseries}{\thesection}{15pt}{\Large}
\titleformat{\subsection}{\large\bfseries}{\thesubsection}{12pt}{\large}

\newenvironment{sidebox}{\begin{framed}\noindent}{\end{framed}}
\newenvironment{doctrine}{\begin{quote}\itshape}{\end{quote}}

\newcommand{\ninth}{\textcolor{everflame}{\textbf{※}}}

\begin{document}

\title{The Two Fires}
\subtitle{Schism, Censure, and the War for the Soul of the Faithful}
\author{An Expansion for Fate's Edge}
\date{Under the lamp, bread and salt --- the flame witnesses both sides}
\maketitle

\tableofcontents

\chapter*{A Note on Method}
\addcontentsline{toc}{chapter}{A Note on Method}

This expansion was written in a candle-lit scriptorium in Thepyrgos, on paper that had been used to wrap heretical pamphlets. The ink is lamp-black and the ash of a burned confession. The author is a convert who has worn both collars and will not say which fire she chose.

The Everflame and the Temple of Light share the same Patrons. They recite similar prayers. They light similar candles.

They are not the same faith.

One is the state religion of an empire. One is the hedge-court faith of a fractured kingdom. One burns heretics at the stake. One burns them in effigy and calls it rehabilitation.

This expansion is not a polemic. It is a ledger. Both sides owe debts. Both sides have paid in blood. The reader is not asked to choose. The reader is asked to witness.

\citation{--- Sister Mariam, Order of Mercy and Grace (defrocked), Thepyrgos, year of the schism}

\part{The Two Fires}

\chapter{Working Title and Core Concept}

\section{The Schism in Brief}

The Everflame and the Temple of Light were once the same faith. The rupture occurred during the Utaran succession crisis, when the High Auditor of Ecktoria declared that the flame could not be carried across the mountains without losing its purity. The Synod of Thepyrgos responded that purity was not the point---illumination was.

Three centuries later, neither side remembers the original argument. They remember the burnings.

\section{What Both Faiths Share}

\begin{itemize}
\item The same three Patrons: \patron{Oath of Flame \& Light}, \patron{Inquisitor Prime}, \patron{Oath of Mercy and Grace}.
\item The same sacred texts (though the Everflame uses a redacted edition and the Temple uses a glossed edition).
\item The same prayers (recited in different orders, with different emphases).
\item The same enmities (Malachai cults, Daughters of the Cord, unsanctioned magic).
\end{itemize}

\section{What Divides Them}

\begin{tabularx}{\textwidth}{X X X}
\toprule
\textbf{Aspect} & \textbf{Everflame (Ecktoria)} & \textbf{Temple of Light (Viterra/Thepyrgos)} \\
\midrule
Nature of the divine & Flame: consuming, purifying, juridical & Light: illuminating, refracting, pastoral \\
Relation to state & Imperial religion & Civic religion with exemptions \\
Heresy & Crime against the empire & Crime against the soul \\
Treatment of heretics & Execution (stake) & Rehabilitation (re-education) \\
Attitude to magic & Unsanctioned magic = heresy & Unsanctioned magic = mental illness \\
Inquisitorial arm & Flame-Wardens & Chain-Lanterns \\
Attitude to children with talent & Burn the taint & Conscript the child \\
\bottomrule
\end{tabularx}

\section{What This Expansion Adds}

\begin{itemize}
\item Doctrinal depth for both faiths, making them playable and morally complex.
\item The Chain-Lanterns as a full faction, including their conscription of children with arcane talent.
\item New Talents for characters sworn to either Order (or caught between them).
\item The Ninth Citation---a forbidden passage that both faiths suppressed and heretics whisper.
\item A campaign frame set in a border region where the Two Fires compete for souls.
\end{itemize}

\citation{The flame purifies. The light reveals. Neither asks permission. --- Parable of the Two Fires}

\part{The Everflame}

\chapter{History: The Imperial Fire}

\section{Origins in the Utaran Forge}

The Everflame began as a cult of the imperial forge---a belief that the fire which shaped Utaran steel also shaped the empire's soul. When the empire fell, the cult survived, carried by soldiers who had prayed before battle and by smiths who had lit candles for safe hammer-strikes.

Ecktoria adopted the Everflame as its state religion during the Consolidation Wars. The Grand Magistrate's argument was simple: \doctrine{A kingdom without a single flame is a kingdom of shadows.}

\section{The Great Purge (Year 312 AR)}

The first systematic suppression of unsanctioned magic. The Flame-Wardens were established. Hundreds were burned. The \textit{Codex of Illumination} was redacted to remove references to the Ninth Citation. The Temple of Light, then still a reform movement within the Everflame, protested. The protest was ignored.

\section{The Schism (Year 478 AR)}

The Synod of Thepyrgos declared independence, citing the Everflame's corruption by imperial politics. The Grand Magistrate declared the Synod heretical. Both sides excommunicated each other. Neither side has rescinded the excommunication.

\citation{They call us heretics. We call them idolaters. The dead call neither. They are simply dead. --- Ember Prefect manifesto}

\chapter{Doctrine: Truth and Measure Are One}

\section{The Flame as Witness}

The Everflame teaches that fire remembers. A candle lit in testimony cannot be extinguished until the truth is spoken. A heretic burned at the stake does not die---they become \textit{ash-witnesses}, their testimony preserved in the smoke.

\doctrine{The flame does not judge. The flame witnesses. The judge is the empire.}

\section{Confession and Censure}

Confession must be public. The priest lights a candle. The penitent speaks. The candle burns. If the candle gutters before the confession ends, the confession is incomplete---the penitent must return.

Censure is public and permanent. A heretic's name is struck from the parish registry. Their family may not pray for them. Their ashes are scattered beyond the city walls.

\section{The Three Orders}

\begin{itemize}
\item \order{Oath of Flame \& Light}: Warriors and crusaders. The flame is a weapon.
\item \order{Inquisitor Prime (Flame-Wardens)}: Prosecutors and executioners. The flame purges.
\item \order{Oath of Mercy and Grace (Hospitaliers)}: Healers and confessors. The flame comforts.
\end{itemize}

\citation{The flame that does not burn is not a flame. It is a painted candle. --- Grand Magistrate Aurelian IV}

\chapter{Hierarchy: The Grand Magistrate's Seal}

\section{The Grand Magistrate}

Elected by the High Synod of Ecktoria. Serves for life or until removed for heresy. The current Grand Magistrate, His Eminence Varrus the Steadfast, has held the seal for forty-two years.

\section{The Flame-Wardens}

The inquisitorial arm. They wear iron masks and carry censers that smoke with blessed coal. Their authority extends to any territory claimed by Ecktoria. They answer only to the Grand Magistrate.

\section{The Censorate}

Bureaucrats who review confessions, maintain heretical registries, and approve burnings. A Censor can delay an execution indefinitely---or expedite it. They are the most feared office in the Everflame because they are the least visible.

\citation{The Censorate does not burn. The Censorate files. The burning is merely a clerical formality. --- Former Censor, writing anonymously}

\chapter{Practices: The Candle and the Stake}

\section{The Candle Test}

A suspected heretic is seated in a room with a single candle. A priest asks questions. The candle's behavior is observed: steady flame = truth; guttering = deception; extinguishing = demonic possession. The test has no empirical validity. Its results have sent hundreds to the stake.

\section{The Unspoken Ninth Citation}

A procedural loophole. If a heretic's confession cannot be obtained, the Flame-Warden may cite the \textit{Ninth Citation}---a redacted passage that permits execution without confession. The citation is never spoken aloud. It is merely \textit{noted in the ledger}.

\citation{We do not burn the innocent. We burn those whose flames do not match our own. --- Flame-Warden proverb}

\chapter{Relationship to Patrons}

\section{Oath of Flame \& Light (Orthodox)}

The warrior's path. Everflame adherents emphasize the flame's consuming nature. They pray for the strength to burn corruption. Their Rites grant +1 die against heretics and +1 Harm against unsanctioned magic users.

\section{Inquisitor Prime (Prosecutors)}

The judge's path. Everflame inquisitors view the Inquisitor Prime as a divine prosecutor. Their Rites grant +1 die to detect lies and +1 Effect when sentencing the guilty.

\section{Oath of Mercy and Grace (Hospitaliers, Watched)}

The healer's path---but watched. Everflame healers must report any heretic they treat. A confessor who shields a heretic from justice is herself a heretic. The Order of Mercy and Grace has produced more martyrs than any other.

\citation{We heal the body. The Flame-Wardens judge the soul. We are not supposed to notice that the soul is still in the body. --- Hospitalier, speaking off the record}

\part{The Temple of Light}

\chapter{History: The Hedge-Law Faith}

\section{Born from the Margins}

The Temple of Light began as a reform movement within the Everflame---priests who argued that the flame had become a tool of imperial control rather than a source of illumination. They were excommunicated. They did not recant.

\section{The Viterran Refuge}

Viterra's hedge-law kingdom offered sanctuary. The Temple established a synod, built schools, and began translating the \textit{Codex of Illumination} into the common tongue. The faith spread through evangelism, not conquest.

\section{The Chain-Lanterns' Founding}

As the Temple grew, so did the threat of unsanctioned magic. The Synod authorized an inquisitorial order---the Chain-Lanterns---to \textit{investigate and rehabilitate}. Unlike the Everflame's Flame-Wardens, the Chain-Lanterns were not permitted to execute. They found loopholes.

\citation{The Chain-Lanterns do not burn. They conscript. The child who cannot control her magic becomes the hunter who tracks others like her. This is not cruelty. This is efficiency. --- Synod report, classified}

\chapter{Doctrine: Light as Lens}

\section{The Refraction of the Divine}

The Temple teaches that the divine is not a single flame but a spectrum. It refracts into three aspects:

\begin{itemize}
\item \textbf{Gold (Justice)}: The light that exposes wrongdoing.
\item \textbf{Silver (Truth)}: The light that reveals what is hidden.
\item \textbf{White (Mercy)}: The light that forgives.
\end{itemize}

No single aspect is complete without the others.

\section{Confession and Rehabilitation}

Confession is private. The penitent speaks to a priest in a candle-lit room. The priest does not report what is heard. The goal is not punishment but \textit{reintegration}.

Heretics are not burned. They are re-educated. The Temple maintains sanctuaries where the possessed and the wayward are taught to control their gifts. Some leave. Most do not.

\citation{We do not burn the tainted. We quarantine them. We do not execute the mad. We medicate them. The difference is not as large as we pretend. --- Chain-Lantern confessor}

\chapter{Hierarchy: The Synod and the Chain}

\section{The Synod}

A council of bishops, abbots, and lay scholars. Meets twice yearly in Thepyrgos. Elects the Archon of Light, who serves as spiritual leader but has no temporal authority.

\section{The Dawn-Knights}

Warriors who protect pilgrims, guard sanctuaries, and---when necessary---hunt rogue magic users. Unlike the Everflame's crusaders, the Dawn-Knights are forbidden from initiating violence. They may only respond.

\section{The Chain-Lanterns}

The Temple's inquisitorial arm. They carry iron lanterns on chains---the flame is visible, the chain is a reminder of restraint. They investigate, detain, and rehabilitate. They are also the Temple's spies, assassins, and secret police.

\citation{The Chain-Lanterns are the Temple's shadow. The Dawn-Knights are its sword. The Synod pretends neither exists. --- The Gray Wanderer}

\chapter{Practices: The Bell and the Mirror}

\section{The Bell-Toll Ordeal}

A suspect is placed in a chamber with a single bell. They must answer questions. Each lie causes the bell to ring---or so the suspect is told. The bell is rung by a hidden Chain-Lantern. The ordeal has no empirical validity. Its results have sent hundreds to re-education.

\section{The White Mirror}

A polished silver surface used to detect unsanctioned magic. A mage who looks into the mirror sees their own reflection---but the reflection is slightly delayed. The longer the delay, the more dangerous the mage. The mirror is not magical. The mirror is a prop. The Chain-Lanterns' judgment is the only test that matters.

\citation{The white mirror shows what the Chain-Lantern wants to see. That is not cynicism. That is policy. --- Former Chain-Lantern, testifying anonymously}

\chapter{Relationship to Patrons}

\section{Oath of Flame \& Light (Dawn-Knights)}

The warrior's path---but the Temple emphasizes protection, not destruction. Dawn-Knights pray for the strength to shield the innocent. Their Rites grant +1 die when defending another and +1 Armor when protecting a sanctuary.

\section{Inquisitor Prime (Chain-Lanterns)}

The judge's path---but the Temple emphasizes investigation, not sentencing. Chain-Lanterns pray for the patience to find the truth. Their Rites grant +1 die to detect lies and +1 Effect when exposing corruption.

\section{Oath of Mercy and Grace (Hospitallers)}

The healer's path---and the Temple gives them autonomy. Hospitallers may treat anyone, regardless of faith or crime. They are not required to report what they hear. This autonomy has produced the Temple's most beloved saints---and its most dangerous secrets.

\citation{I treated a heretic. I did not report him. He became my teacher. I am still a Hospitalier. I am also a heretic. The oath does not forbid contradiction. It only forbids cruelty. --- Sister Mariam, defrocked}

\part{The Chain-Lanterns}

\chapter{The Temple's Sword and Shadow}

\section{History of the Order}

Founded during the Viterran witch-panics of the seventh century. The Synod needed an inquisitorial arm that could investigate without provoking the hedge-courts. The Chain-Lanterns were that arm.

Initially, they were respected. They exposed genuine conspiracies and saved innocents from mob justice. Over time, their methods grew more aggressive. The chain that symbolized restraint became a weapon.

\section{Methods and Tools}

\begin{itemize}
\item \textbf{The Candle-Smoke Test}: A suspect breathes over a candle. The smoke's color indicates guilt---or so the Chain-Lanterns claim. The test's results are predetermined.
\item \textbf{The Bell-Toll Ordeal}: Described above. The Chain-Lanterns have perfected the timing.
\item \textbf{The White Mirror}: Described above. A prop that has never failed to produce a confession.
\end{itemize}

\section{Conscription of Children}

The Temple's most controversial practice. Children with arcane talent are identified---often through the Chain-Lanterns' network of informants---and taken to sanctuaries for "rehabilitation." Most never leave. Some become Chain-Lanterns themselves.

\citation{They saved me from the fire. Then they taught me to light it. I am grateful. I am also a monster. The two are not exclusive. --- Former Chain-Lantern child conscript}

\section{Secret Operations in Ecktoria}

The Chain-Lanterns operate illegally in Everflame territory, hunting "heretics" who have fled Temple justice. The Everflame's Flame-Wardens know. They have cooperated---and clashed. The Chain-Lanterns have a safehouse in Ecktoria's catacombs. Its location is a secret the Everflame would kill to learn.

\citation{They call us heretics. We call them idolaters. We both call each other at midnight, sometimes, to share information about Malachai cults. The enemy of my enemy is still my enemy---but we can be polite about it. --- Chain-Lantern and Flame-Warden, meeting notes}

\part{The Shared Patrons}

\chapter{Oath of Flame \& Light: The Warrior's Path}

\section{Everflame Interpretation: The Flame Is a Weapon}

The Everflame's warriors see themselves as instruments of divine judgment. They are crusaders. They do not negotiate with heresy. They burn it.

\section{Temple Interpretation: The Light Is a Guide}

The Temple's Dawn-Knights see themselves as shields. They protect the innocent. They do not seek out heresy---they wait for it to reveal itself.

\section{The Schism Within}

The shared Order is riven by the Two Fires. A Dawn-Knight who refuses to burn a heretic is a coward to the Everflame. A crusader who burns a village is a monster to the Temple. Neither side is entirely wrong.

\section{New Talents}

\subsection{Radiant Vow (Everflame, 6 XP)}

Swear an oath before a flame. While keeping the oath, gain +1 die to all actions. Breaking the oath causes 2 Harm (the flame's rebuke). The oath's terms must be specific and achievable.

\subsection{Dawn's Shield (Temple, 6 XP)}

Once per session, interpose yourself to protect an innocent. Convert any Harm they would have taken into Fatigue for yourself. The innocence must be genuine---the shield does not protect the guilty.

\subsection{Candle-Knight (Hybrid, 10 XP)}

You have studied both traditions. Once per session, you may invoke either faith's interpretation of a shared Rite, but you gain a \texttt{Split Loyalty [6]} timer. When it fills, you must choose a side. Refusal is heresy to both.

\citation{I have burned heretics. I have shielded the accused. I do not know which fire is mine. I light both candles. --- Candle-Knight, speaking at a crossroads}

\chapter{Inquisitor Prime: The Judge's Path}

\section{Everflame Interpretation: The Flame-Wardens}

Prosecutors who burn corruption. They are not investigators---they are executioners. Their authority is absolute within Everflame territory.

\section{Temple Interpretation: The Chain-Lanterns}

Investigators who expose heresy. They are not executioners---they detain and rehabilitate. Their authority is constrained by Synod oversight.

\section{The Schism Within}

The shared Order is a cold war. Flame-Wardens view Chain-Lanterns as fanatics who have forgotten that mercy is a virtue. Chain-Lanterns view Flame-Wardens as state-sponsored murderers. They have cooperated against common enemies. They have also framed each other for heresy.

\section{New Talents}

\subsection{Flame-Warden's Brand (Everflame, 8 XP)}

Mark a heretic. For one session, you know their location (as long as they remain within the same region). They suffer --1 die to resist your social rolls. The brand is visible to other Flame-Wardens.

\subsection{Chain-Lantern's Bell (Temple, 8 XP)}

Ring a bell. For one scene, any unsanctioned magic within Near must test Spirit+Resolve (DV 4) or be revealed---the caster's location, the spell's nature, and the caster's face. The bell's note carries. Witnesses will know something happened.

\subsection{The Ninth Citation (Heresy, 12 XP)}

You have read the forbidden passage. Once per session, you may invoke it to gain +2 dice on any roll---but both faiths will hunt you. The citation cannot be used twice in the same location. The Unspoken notices repetition.

\citation{The Ninth Citation is not a spell. It is a door. I have read it. I have not opened the door. I am not brave enough. --- Anonymous heretic}

\chapter{Oath of Mercy and Grace: The Healer's Path}

\section{Everflame Interpretation: Mercy Is Earned}

The Everflame's Hospitaliers heal heretics---but they must report them. A heretic who receives care and does not confess is returned to the Flame-Wardens. The healing is not free. The debt must be paid.

\section{Temple Interpretation: Mercy Is Given}

The Temple's Hospitallers heal anyone, regardless of faith or crime. They are not required to report. The healing is free. The only debt is gratitude---and gratitude is not tracked in ledgers.

\section{The Schism Within}

The shared Order is a battlefield of conscience. Everflame Hospitaliers who shield heretics are martyrs. Temple Hospitallers who report heretics are traitors. The same action, judged by different fires, produces opposite verdicts.

\section{New Talents}

\subsection{Confessional Seal (Everflame, 4 XP)}

You cannot be magically compelled to reveal a confession. You also cannot break a confidence voluntarily without marking 1 Fatigue. The seal is absolute. Break it, and the flame will remember.

\subsection{Unburdened Healing (Temple, 4 XP)}

When you heal someone, you may choose not to report them (even if your faith requires it). Mark 1 Fatigue if you conceal a heretic. The fatigue is the weight of the secret.

\subsection{The Ninth Prayer (Heresy, 12 XP)}

You have learned to heal without conditions---no confession, no reporting, no debt. Each time you use this talent, mark 1 Corruption. The Unspoken is watching. The Unspoken is pleased.

\citation{They taught me to heal. Then they taught me to judge. I rejected the judgment. I kept the healing. Both faiths consider me a heretic. The sick do not care. --- The Ninth Prayer, first known practitioner}

\part{The Quiet War}

\chapter{The Battle for the Mistlands}

Both faiths compete for souls in the Mistlands. The Everflame sends missionaries with guards. The Temple sends Dawn-Knights disguised as pilgrims. The locals accept the medicine and burn the pamphlets.

\section{New Faction: The Gray Pilgrims}

Neutral healers who accept aid from both faiths and convert to neither. They are respected by the Mistlanders and distrusted by both churches. A Gray Pilgrim who accepts healing from the Everflame must also accept a lecture on sin. A Gray Pilgrim who accepts grain from the Temple must also accept a child's catechism.

\section{New Threat: The Chain-Lanterns "Investigating"}

The Chain-Lanterns have begun operating in the Mistlands, hunting witches who may be midwives with herb gardens. The Everflame views this as an incursion. The Mistlanders view it as terrorism. The Chain-Lanterns view it as justice.

\citation{The Mistlands do not belong to the Everflame. They do not belong to the Temple. They belong to the fog. The fog does not care about your schism. --- Mistlander elder}

\chapter{The Linn Crusade}

The Temple sends no armies to Linn---only teachers, healers, and grain. The Linn accept the grain. The children learn to read the \textit{Codex of Illumination}. The old shrines fall silent.

\section{New Faction: The Shrine-Devout}

Linn who resist conversion. They keep the old shrines, light the old candles, and whisper the old names. The Temple views them as obstacles. The Everflame views them as pagans. The Shrine-Devout view themselves as the last witnesses.

\section{New Mechanic: Faith Pressure [6]}

A regional timer that ticks as the Temple gains influence. Advances when the Temple builds a school, converts a chieftain, or baptizes a child. Retreats when the Shrine-Devout reclaim a sacred site, expose a corrupt missionary, or drive out a Chain-Lantern cell. When full, the region converts (or rebels). The GM decides which.

\citation{They built a school. They taught my daughter to read. She read the \textit{Codex}. She asked me why our ancestors burned candles. I could not answer. The old shrines fell silent. --- Linn father, converted}

\chapter{Ecktoria's Shadow War}

The Chain-Lanterns operate illegally in Ecktoria, hunting "heretics" who have fled Temple justice. The Everflame's Flame-Wardens view this as a violation of sovereignty. The two inquisitions have fought, betrayed, and occasionally cooperated.

\section{New Faction: The Ember Prefects}

Everflame heretics who believe the flame belongs in the streets, not basilicas. They are hunted by both the Everflame (as schismatics) and the Temple (as heretics). They may be the only honest actors in the Quiet War.

\section{New Threat: The Catacomb Cell}

A Chain-Lantern cell operating in Ecktoria's catacombs, supported by unknown patrons. The cell has evaded capture for years. Its members know the Ninth Citation. Its leader has never been seen without a white mirror.

\citation{We are not spies. We are witnesses. The Everflame has forgotten what the flame is for. We are here to remind them. --- Chain-Lantern cell leader, recorded confession}

\chapter{Thepyrgos: The City of Two Fires}

Thepyrgos is the heart of the Temple of Light---but the Everflame has a presence, and the Covenant has a foothold. The Chain-Lanterns conscript children with arcane talent. The Everflame's Flame-Wardens investigate. The Synod tries to keep the peace.

\section{New Faction: The Synod}

The Temple's leadership. They are scholars, not warriors. They debate theology while the Chain-Lanterns conscript children. They are not corrupt. They are not naive. They are caught between their principles and their survival.

\section{New Threat: The Ember Prefect Cell}

Everflame heretics planting evidence to trigger a purge. They want the Chain-Lanterns expelled from Ecktoria. They are willing to burn innocents to achieve that goal. The party must stop them---or help them.

\citation{The Synod debates. The Chain-Lanterns act. The children vanish. This is not a criticism. This is a description. --- The Gray Wanderer, Thepyrgos}

\part{The Covenant of Mykkiel}

\chapter{History: The Faith of the Contract}

\section{Origins in the Pereshi Highlands}

The Covenant did not begin in a temple or a court. It began in a marketplace, where a Pereshi merchant and a Sidhi trader could not agree on the price of salt. A judge was called—not a priest, not a warlord, but a witness who knew the weight of a promise. That judge wrote down what was said, and both parties signed. The salt was sold. The debt was recorded. The Covenant was born.

The Pereshi have always been a people of ledgers and caravans, and their law preceded their gods. When the Utaran Empire expanded east, it found the Pereshi already bound by a code more ancient than imperial edict. The Empire did not suppress it—it adopted it, codified it, and called it the \textit{Canons of Obligation}. When the Empire fell, the Covenant survived, carried by notaries, judges, and the weight of the written word.

\section{The Covenant and the Fall of Utar}

The Utaran Empire's collapse was slow, but the Covenant was not caught in the rubble. Its judges had always served both emperor and merchant; when the emperor vanished, they continued to serve the contract. The Covenant's courts remained open in Galanina, in the Fhara trade routes, and in the secret cells of Ashaan. The Empire was a client, not a patron. The Covenant outlasted it.

\citation{The Empire fell. The Covenant did not. The Empire was a fire. The Covenant is a ledger. Fire burns out. Ledgers are forever. --- Pereshi proverb}

\section{The Ashaan Schism}

In Ashaan, the Covenant's role is more contentious. The Veiled Sisters, the sorcerer-queens who once ruled the Pereshi highlands, were oath-breakers by Covenant law—they had seized power without contract, ruled without consent, and bound spirits without witnesses. The Covenant's judges condemned them \textit{in absentia}, invoking the Ninth Clause to try them without presence. When the Ashaan Crusade liberated Galanina, the Covenant's courts were the first to reopen. They have not closed since.

But the Covenant is not a religion of conquest. It does not seek converts; it seeks \textit{signatories}. In Ashaan, the Covenant's authority is accepted because it is useful—a neutral arbiter in a land of rival magocracies. The Veiled Sisters are gone, but their legacy of broken oaths lingers. The Covenant's judges are the ones who decide which debts remain unpaid.

\citation{We do not forgive. We do not forget. We record. The ledger is eternal. --- Covenant judge, Galanina}

\chapter{Theology: The Law as Witness}

\section{The Divine Contract}

The Covenant teaches that the universe is not a creation but an \textit{agreement}. In the beginning, the great powers—whom outsiders call Patrons—entered into a binding contract to govern the chaos of existence. That contract is the Covenant, and it endures not because it is enforced but because it is \textit{true}. Those who honor the Covenant align themselves with the fundamental structure of reality; those who break it become dissonant, unmoored, and eventually erased by the world's rejection.

\doctrine{The universe does not punish oath-breakers. The universe merely remembers their debt. The debt is the punishment.}

\section{The Role of Mykkiel}

Mykkiel, called the Arbiter, is not the \textit{source} of the Covenant but its \textit{living seal}. He is the Patron who witnesses every oath and records every breach. He does not create law; he \textit{remembers} it. His followers are judges, not legislators. They do not interpret the Covenant; they \textit{apply} it.

\section{The Grand Ledger}

Every judgment of a Covenant court is recorded in the Grand Ledger, a brass-bound book kept in a vault in Galanina. The Ledger is said to weigh more than a horse. A name struck from the Ledger is struck from the Covenant community—they are no longer bound by the Covenant's protections and may be treated as outlaws. The Ledger is the physical manifestation of the Covenant's authority.

\citation{The Grand Ledger is not a book. It is a contract between the living and the dead. A name written in it is not merely recorded; it is \textit{bound}. --- Keeper of the Ledger}

\section{The Unspoken Ninth Clause}

The Covenant has its own forbidden passage, analogous to the Everflame's Ninth Citation. The \textit{Unspoken Ninth Clause} is a ruling that was never recorded—a judgment that exists only in the memory of the judge who pronounced it. The ruling holds that "in cases where the accused is not present, and no witness can be found, the judge may pronounce judgment based on the weight of testimony already recorded."

The Ninth Clause is used to justify trials \textit{in absentia}—condemning the dead, the exiled, and the unreachable. It is also used to justify \textit{preemptive judgment}—convicting someone before they have committed the crime, on the grounds that "the pattern of their oath-breaking is already evident."

The Covenant's orthodox authorities deny that the Unspoken Ninth Clause exists. They know that it does. The clause is a secret passed from judge to judge, never written down, but always remembered.

\citation{The Ninth Clause is not a loophole. It is a mirror. It shows us what we already know. It reminds us that the law is not a cage; it is a lens. --- Circuit Judge, speaking off the record}

\chapter{The Three Orders of the Covenant}

\section{The Order of Mykkiel (Orthodox)}

The judges, notaries, and law-scholars. They settle disputes, draft contracts, and administer the Covenant Courts. Their authority is recognized wherever the Covenant has a foothold—which is to say, wherever trade and treaties require a trusted third party.

\begin{itemize}
\item \textbf{Rites}: Those of binding oaths, weighing testimony, and sealing judgments. The Order's Rites are slow, deliberate, and \textit{final}. A judgment pronounced in a Covenant Court is not appealed; it is \textit{enforced}.
\item \textbf{Organization}: Hierarchical, with High Justiciars at the top (elected by the Covenant Council) and Circuit Judges who travel the roads. Each court has a Seal-Bearer, a Notary, and a Clerk of Records.
\item \textbf{Where They Are Found}: Galanina, the Fhara trade routes, the secret cells of Ashaan, and increasingly in Thepyrgos. In Ecktoria, they are viewed as heretics. In Vhasia, they are tolerated as a philosophy.
\end{itemize}

\section{The Order of the Inquisitor Prime (Prosecutorial)}

In the Covenant, the Inquisitor Prime is the \textit{Accuser}—the one who brings hidden sins to light. They are not hunters of heretics but \textit{auditors of the soul}. They investigate broken oaths, expose secret perjury, and demand restitution. Their Rites are those of revelation and correction. They carry no swords—only iron scales and a copy of the \textit{Canons of Obligation}.

\begin{itemize}
\item \textbf{Rites}: Those of investigation, interrogation, and the weighing of testimony. They do not execute; they \textit{expose}.
\item \textbf{Organization}: Secretive and decentralized, with each cell operating independently. Identifiable by their iron scales and violet sash.
\item \textbf{Where They Are Found}: Galanina, the Fhara, and as resistance fighters in Ashaan. In Thepyrgos, they are feared as spies.
\end{itemize}

\section{The Order of the Oath of Flame \& Light (Hospitalier—Healing Focus)}

In the Covenant, the Oath of Flame \& Light is not a warrior's path but a \textit{healer's}. They tend to the wounded in body and soul, but their healing always comes with a \textit{contract}: the patient must swear to live a more lawful life. Their Rites are those of purification, cleansing, and bound mercy.

\begin{itemize}
\item \textbf{Rites}: Healing, purification, and the binding of mercy. A healer will not treat a murderer unless the murderer agrees to face judgment; they will not comfort a liar unless the liar agrees to confess.
\item \textbf{Organization}: Small and tightly knit, with each healer attached to a specific court or judge. Sworn to secrecy regarding confessions but required to report \textit{knowledge of unrepented crimes}.
\item \textbf{Where They Are Found}: Wherever the Covenant has a foothold, including refugee camps in Thepyrgos.
\end{itemize}

\citation{The Three Orders are three aspects of the same law. The judge pronounces. The accuser exposes. The healer mends. None can stand without the others. --- Covenant creed}

\chapter{The Covenant and the Two Fires}

\section{Areas of Cooperation}

The Covenant maintains a wary distance from both the Everflame and the Temple of Light. It regards the Everflame as \textit{juridically imperial}—the Everflame's laws are the laws of the empire, not the laws of the Covenant. It regards the Temple of Light as \textit{pastorally lax}—the Temple's emphasis on mercy and rehabilitation is a corruption of the divine law.

Yet the Covenant shares certain principles with both faiths. Like the Everflame, it believes that \textit{law is sacred}. Like the Temple, it believes that \textit{witness is binding}. The Covenant's judges often cooperate with Everflame magistrates and Temple clergy, despite theological disagreements. A case involving a broken oath may be heard in a Covenant Court, regardless of the oath's origin.

\begin{itemize}
\item \textbf{Shared Inquisitions}: When a heretic or oath-breaker flees across borders, the Covenant's Inquisitors, the Everflame's Flame-Wardens, and the Temple's Chain-Lanterns may coordinate pursuit. The cooperation is transactional—each side provides information to the others—but it is effective.
\item \textbf{Defense of the Faithful}: When a community of Covenant followers is threatened (by Malachai cults, Daughters of the Cord, or unsanctioned magic), the Covenant may call on the Everflame's military or the Temple's Dawn-Knights for protection. The Call of the Covenant is rarely refused; to refuse is to admit that one does not honor the law.
\item \textbf{Crusades}: When a region becomes tainted by unsanctioned magic or demonic influence, the Covenant may call for a Crusade—a unified military campaign involving all three faiths.
\end{itemize}

\section{Areas of Conflict}

\begin{itemize}
\item \textbf{Jurisdiction}: The Covenant claims universal authority over all oaths and contracts. The Everflame claims universal authority over all heresy. The Temple claims universal authority over all miracles. When a case involves all three, no one can agree on who has jurisdiction.
\item \textbf{The Chain-Lanterns}: The Temple's inquisitors are the Covenant's most frequent opponents. Chain-Lanterns routinely ignore Covenant jurisdiction, arguing that "the Temple's authority supersedes all worldly law." Covenant judges have issued warrants for the arrest of Chain-Lanterns, but the warrants are never served.
\item \textbf{The Grand Ledger}: The Everflame's Grand Magistrate has demanded access to the Covenant's Grand Ledger, arguing that "the empire's law supersedes the Covenant's." The Covenant has refused. The dispute has not been resolved.
\end{itemize}

\citation{The Two Fires burn bright. The Covenant burns slow. One is a flame. One is a ledger. Both are necessary. Neither is sufficient. --- Pereshi philosopher}

\chapter{The Covenant and the Crusades}

\section{The Ashaan Crusade (890–901 AR)}

The Ashaan Crusade was the largest military expedition in modern memory. Three faiths—the Everflame, the Temple of Light, and the Covenant—joined forces to liberate the Pereshi highlands from the Veiled Sisters, a witch-queen's regime that had ruled for three centuries.

The Covenant was the driving force behind the Crusade. Its judges had long memories of the Veiled Sisters' crimes; they saw the Crusade as a divine audit. After the Crusade, the Covenant established a permanent court in Galanina, which still operates today. The court has issued more than a thousand judgments against former Veiled Sisters, their allies, and their descendants.

The Crusade succeeded, but the coalition fractured almost immediately. The Everflame claimed imperial authority over the liberated territories; the Temple claimed pastoral authority; the Covenant claimed legal authority. A small war broke out among the Crusaders. The people of Galanina, tired of all three, elected a Pereshi council and declared independence.

\citation{We fought side by side. We won. Then we fought each other. The victory was not ours. It belonged to the people we forgot to include. --- Covenant judge, reflecting on the Crusade}

\section{The Linn Crusade (Current)}

The Linn Crusade is not a military campaign—it is a Crusade of education, conversion, and cultural replacement. The Temple of Light sends teachers, healers, and grain to Linn. The Everflame sends missionaries, guards, and basalt basilicas. The Covenant sends judges and notaries, offering to mediate disputes between the two.

The Linn welcome the gifts and convert to neither. The Covenant's role is neutral arbitration—a third party that both faiths distrust but both need.

\citation{In Linn, the Temple builds schools. The Everflame builds basilicas. The Covenant builds courthouses. The Linn build bonfires and burn the pamphlets. --- The Gray Wanderer}

\section{The Covenant's Crusade Doctrine}

The Covenant views Crusades as \textit{properly ordered violence}—a form of law enforcement on a divine scale. A Crusade is not a war of conquest but a \textit{judgment}: a verdict pronounced by the divine, executed by the faithful.

\begin{itemize}
\item \textbf{Judges and Notaries}: The Covenant supplies judges, notaries, and legal counsel to Crusader armies. A Crusader who commits a crime is tried in a Covenant court; a Crusader who makes an oath is bound by Covenant law. The Covenant's presence is intended to ensure that the Crusade remains just.
\item \textbf{The Grand Ledger}: The Covenant maintains a separate ledger for Crusades, recording every oath made and every judgment pronounced. This ledger is used to settle disputes after the Crusade ends—and to pursue those who broke their oaths.
\item \textbf{The Ninth Clause}: In extreme cases, the Covenant may invoke the Ninth Clause to try enemies \textit{in absentia}. This allows Crusader commanders to pronounce judgment on absent foes, condemning them to excommunication or outlawry without a trial.
\end{itemize}

\citation{A Crusade is not a war. It is an audit. The sword is a quill. The blood is ink. The ledger is the only thing that matters. --- Covenant Crusade doctrine}

\chapter{New Talents and Rites for the Covenant}

\section{Talents}

\subsection{Covenant Justiciar (Prestige, 10 XP)}

\emph{Prerequisites: Spirit 3+, Presence 2+, Rite of the Consecrated Mandate or Cantor of Mykkiel.}

\textbf{Effect:} You are a sworn judge of the Covenant. Once per session, you may declare a single oath you have sworn (to protect a person, to destroy a foe, to uphold a law) as inviolable. For the remainder of the session, you gain +2 dice to all rolls that directly advance that oath, and you ignore the first Harm 2 each scene that would prevent you from fulfilling it.

\textbf{Cost:} After the scene, you gain a permanent Debt Scar (The Unbroken Word): your voice carries a faint echo of the Covenant, and you can never again speak a false oath without the scar flaring painfully (1 Fatigue). You may ignore one debt or obligation that conflicts with your sworn oath, but each future oath you swear increases your Obligation to Mykkiel by 1 permanently.

\subsection{Seal of the Grand Ledger (Minor, 4 XP)}

\emph{Prerequisites: Presence 2+, Lore 1+.}

\textbf{Effect:} You carry a small iron seal that, when pressed onto a document or threshold, makes the oath sworn there binding. Any creature that knowingly breaks an oath sealed by you suffers -1 die to all rolls until they make amends. You may use this seal once per session.

\subsection{Auditor's Gaze (Major, 6 XP)}

\emph{Prerequisites: Wits 3+, Insight 2+.}

\textbf{Effect:} Once per scene, when you observe a transaction, negotiation, or argument, you may ask the GM one question about a hidden obligation, a broken promise, or a concealed motivation. The GM must answer truthfully. On a partial success (or if the GM spends 1 SB), the target becomes aware they are being audited.

\subsection{Witness of the Unnamed Price (Prestige, 6 XP)}

\emph{Prerequisites: Spirit 3+, Lore 2+.}

\textbf{Effect:} When you mark someone guilty (fictionally), once per session you may shift one failed roll into a partial—but mark 1 Fatigue for the burden of judgment.

\subsection{Mercy as Blade (Minor, 4 XP)}

\emph{Prerequisites: Presence 2+.}

\textbf{Effect:} When you spare someone, gain Position +1 against them until they betray or defy your mercy.

\subsection{Doctrine-Bearer (Minor, 4 XP)}

\emph{Prerequisites: Lore 2+.}

\textbf{Effect:} Once per session, declare one outcome as precedent. Allies gain +1 die when invoking your ruling; enemies do so too in appeal.

\section{Rites of Mykkiel}

\subsection{Rite of the Consecrated Mandate (Basic, 5 XP)}

\textit{Scene; Zone; No.}

\textbf{Materials:} A war-banner sanctified with sacred oils, or sanctified salt and ritual boundaries.

\textbf{Effect:} Consecrate an area for a righteous cause. Followers gain +1 die to resist fear; enemies suffer –1 die to oppose your proclamation. Create a 4-segment Divine Mandate timer.

\textbf{Push It:} The banner emits visible radiance; unbelievers must test Resolve (DV 3) or falter, taking –1 die to attacks for the scene.

\subsection{Rite of the Covenant Seal (Advanced, 9 XP)}

\textit{Extended; Zone; No.}

\textbf{Materials:} Wax, sacred seals, the names of sworn parties.

\textbf{Effect:} Formalize a binding covenant. All participants suffer concrete penalties for breach (–1 die to relevant actions). Creates an 8-segment Covenant Bond timer; each major breach advances the timer by 2. When full, the covenant shatters and all parties mark +1 Obligation.

\subsection{Rite of Divine Judgment (Master, 13 XP)}

\textit{Extended; Near; Yes.}

\textbf{Materials:} Scales of justice, testimony (written or spoken), ritual condemnation.

\textbf{Effect:} Pronounce divine judgment. The target suffers escalating penalties (–1 die, then –2 dice) on all social and contested rolls until they atone. Serious crimes (oath-breaking, sacrilege, murder) may warrant greater effects (e.g., Harm 2, exile, or the attention of Covenant inquisitors). Start a 10-segment Judgment timer; each sunrise advances it by 1. When full, the judgment becomes permanent until a greater authority lifts it.

\textbf{Push It:} The judgment becomes public knowledge; the community shuns an unrepentant target.

\subsection{Rite of the Chosen Legion (Master, 14 XP)}

\textit{Extended; Large Zone; No.}

\textbf{Materials:} Consecrated ground, ritual weapons, an oath of fidelity sworn by each participant.

\textbf{Effect:} Sanctify an area for the faithful. Believers gain +1 die to defense; unbelievers suffer –1 die. Supernatural desecration within the zone risks immediate retaliation: the GM gains 1 SB and a divine warning manifests.

\textbf{Push It:} The area becomes permanently consecrated; the timer instead tracks the arrival of a celestial guardian.

\citation{Mykkiel's Rites are not prayers. They are procedures. They do not ask for mercy. They demand compliance. --- Covenant catechism}

\chapter{Covenant Influence and Obligation}

\section{Covenant Influence Track}

When a character serves the Covenant—as a judge, inquisitor, or healer—they gain \textbf{Covenant Influence}. This is a reputation track that unlocks access to Covenant resources and responsibilities.

\begin{center}
\begin{tabular}{p{3cm} p{5cm} p{6cm}}
\toprule
\textbf{Influence Level} & \textbf{Benefits} & \textbf{Responsibilities} \\
\midrule
0–2 & Access to Covenant courts as a petitioner. & None. \\
3–4 & May request a Circuit Judge's aid; gain +1 die to legal negotiations. & Must render judgment in a dispute when asked. \\
5–6 & Access to the Grand Ledger (one query per session); may issue a Writ of Censure. & Must participate in an Inquisition when summoned. \\
7+ & May convene a Covenant Court; gain +2 dice to Command in legal matters. & Must investigate any reported oath-breaking in their region. \\
\bottomrule
\end{tabular}
\end{center}

\textbf{Clearing Influence:} Influence clears when you refuse a summons, break a Covenant oath, or fail to enforce a judgment. Each violation reduces Influence by 2 and may trigger a penalty (excommunication, loss of Rites, or a Debt Timer).

\citation{Influence is not a reward. It is a leash. The higher you rise, the tighter it pulls. --- Circuit Judge}

\chapter{Adventure Hooks and Campaign Frames}

\section{Adventure Seeds}

\subsection{The Missing Ledger}

A Covenant Circuit Judge has vanished while investigating a series of broken oaths in the Mistlands. The judge's ledger, containing sensitive testimony, is missing. The party must find the judge—or the ledger—before the Chain-Lanterns do.

\subsection{The Ninth Clause}

A Covenant court has invoked the Ninth Clause to condemn an absent noble for heresy. The noble's family begs the party to find evidence of their innocence. The twist: the noble is guilty, but the Ninth Clause was invoked illegally. The party must expose the court's overreach—or prove the noble's guilt beyond doubt.

\subsection{The Ember Prefect's Rebellion}

An Ember Prefect cell in Thepyrgos plans to burn a Temple basilica and blame the Chain-Lanterns. The Covenant's judges have learned of the plot and want the party to stop it—but the judges also want to use the incident to curtail the Chain-Lanterns' authority. The party must choose: prevent the attack, expose the Ember Prefects, or weaponize the plot for the Covenant.

\subsection{The Linn Crusade}

The party is hired as neutral mediators in a Linn village where the Temple and the Everflame are competing for converts. Each side has committed atrocities; the villagers are caught in the middle. The party must broker a peace—or choose a side and face the consequences.

\subsection{The Debt of the Unbroken Cord}

A Daughter of the Cord has been accused of oath-breaking by a Covenant court. The Daughter claims the oath was void because the Cord's authority supersedes the Covenant's. The party must investigate the truth—and decide which faith's law holds sway.

\section{Campaign Frame: The Two Fires War}

A full-length campaign (10–15 sessions) set in a border region where both faiths compete for influence. The party chooses a side—or tries to stay neutral as the fires consume everything.

\subsection{Timers}

\begin{itemize}
\item \textbf{Everflame Influence [8]}: Ticks when the Everflame converts a village, builds a basilica, or burns a heretic. When full, the region is under Everflame control.
\item \textbf{Temple Influence [8]}: Ticks when the Temple builds a school, heals a chieftain, or conscripts a child. When full, the region is under Temple control.
\item \textbf{Civil Unrest [6]}: Ticks when the two faiths clash. When full, the region erupts in open conflict. Both Influence timers pause. Violence is the only language.
\item \textbf{Heresy Hunt [4]}: Ticks when the party uses the Ninth Citation or shelters a heretic. When full, a joint inquisition (Flame-Wardens and Chain-Lanterns) arrives. The party must flee, fight, or confess.
\item \textbf{Covenant Arbitration [6]}: Ticks when the Covenant is invited to mediate. When full, a Covenant judge arrives and imposes a binding settlement. Both sides may reject it—but at the cost of facing the Covenant's judgment.
\end{itemize}

\subsection{Factions}

\begin{itemize}
\item \textbf{Everflame Missionaries}: Zealous, well-funded, and culturally tone-deaf. They offer grain and demand conversion.
\item \textbf{Temple Dawn-Knights}: Polite, persistent, and patient. They offer medicine and ask for catechism.
\item \textbf{Chain-Lanterns (Secret)}: Operating in disguise. They are hunting a specific heretic—or planting evidence to justify intervention.
\item \textbf{Ember Prefects (Heretics)}: Everflame schismatics who believe the flame belongs to the people, not the empire. They may be allies—or terrorists.
\item \textbf{Neutral Locals}: They want the grain, the medicine, and the faiths to leave. They will betray anyone to survive.
\item \textbf{Covenant Arbiters}: They want a binding resolution. They will enforce it with the weight of the Grand Ledger.
\end{itemize}

\subsection{Endings}

\begin{itemize}
\item \textbf{Everflame Victory}: The region converts. Heretics burn. The Temple's sanctuaries are seized. The Chain-Lanterns flee—or die.
\item \textbf{Temple Victory}: The region converts. Unsanctioned magic is suppressed. The Everflame's missionaries are expelled. The Flame-Wardens are hunted.
\item \textbf{Neutral Victory}: Both faiths are expelled. The region declares independence. The party is hailed as liberators—and hunted by both inquisitions.
\item \textbf{Covenant Victory}: The region becomes a Covenant protectorate. Oaths are sworn. Law replaces faith. The Two Fires are extinguished—or subjugated.
\item \textbf{Heresy Victory}: The Ninth Citation is read aloud in a sacred space. The Unspoken answers. The Two Fires are extinguished. Something new begins.
\end{itemize}

\citation{The war does not end when one fire goes out. The war ends when the ledger is balanced. The Covenant's victory is not a flame. It is a signature. --- Covenant judge, final summation}

\chapter*{Final Thought on the Covenant}
\addcontentsline{toc}{chapter}{Final Thought on the Covenant}

The Covenant is not a third fire—it is a ledger. It does not seek to burn or illuminate; it seeks to \textit{record}. Its followers are not zealots or healers; they are auditors, judges, and witnesses. In a world of schism and crusade, the Covenant offers a different path: one of law, duty, and the weight of the word spoken and kept.

The Two Fires will burn until the end of days. The Covenant will outlast them. Because the Covenant is not a fire. It is the \textit{record} of the fire. And records, as the Daughters teach, are forever.

\citation{May your oaths hold true, and your judgments be just. --- The Gray Wanderer, in the shadow of the Grand Ledger}

\appendix

\chapter{Covenant Quick Reference}

\section{Patron's Gift (Mykkiel)}

Once per scene, a Runekeeper of Mykkiel may touch a held item and speak a word of judgment. For the rest of the scene, the item grants +1 die to Investigation or Command rolls when confronting supernatural creatures, and the first time each scene the bearer would be affected by a magical illusion or fear effect, they may reroll the resistance check. After the scene, mark +1 Obligation.

\section{Summary of New Talents}

\begin{tabularx}{\textwidth}{X c c X}
\toprule
\textbf{Talent} & \textbf{Tier} & \textbf{Cost (XP)} & \textbf{Effect} \\
\midrule
Covenant Justiciar & Prestige & 10 & Declare an oath inviolable; +2 dice; ignore first Harm 2 per scene \\
Seal of the Grand Ledger & Minor & 4 & Bind an oath with iron seal; oath-breakers suffer -1 die \\
Auditor's Gaze & Major & 6 & Ask one question about hidden obligations; GM must answer truthfully \\
Witness of the Unnamed Price & Prestige & 6 & Shift a failed roll to partial by marking Fatigue \\
Mercy as Blade & Minor & 4 & Gain Position +1 against those you spare \\
Doctrine-Bearer & Minor & 4 & Declare a precedent; allies gain +1 die when invoking it \\
\bottomrule
\end{tabularx}

\section{Summary of New Rites}

\begin{tabularx}{\textwidth}{X X p{6cm}}
\toprule
\textbf{Rite} & \textbf{Tier} & \textbf{Effect} \\
\midrule
Consecrated Mandate & Basic & Consecrate an area; followers resist fear; enemies suffer -1 die \\
Covenant Seal & Advanced & Formalize a binding covenant; breach creates penalties \\
Divine Judgment & Master & Pronounce judgment; target suffers escalating penalties until atonement \\
Chosen Legion & Master & Sanctify an area; believers gain +1 die to defense; desecration triggers backlash \\
\bottomrule
\end{tabularx}

\citation{These are not spells. They are procedures. Follow them exactly, and the law will hold. --- Covenant practice manual}

\part{Playing Between the Fires}

\chapter{Character Options}

\section{The Convert}

Left one faith for the other. Hunted by your old Order. Distrusted by your new one. You carry the guilt of the fire you abandoned and the coldness of the light you cannot feel.

\section{The Double Agent}

Pledged to both. The weight will break you. You attend the Everflame's masses and the Temple's prayers. You have two confessions. Neither is complete. Both are true.

\section{The Refugee}

Fled the Chain-Lanterns. The Everflame offers sanctuary---at a price. You will be watched. You will be questioned. You will never be trusted. But you will not be conscripted.

\section{The Heretic}

You read the Ninth Citation. You are not supposed to know it exists. Both faiths would kill you if they knew. You carry the forbidden passage in your memory. You have not spoken it aloud. You are not sure you can resist.

\section{The Arbitrator}

You serve the Covenant, the faith of the contract. You are not a convert; you are a signatory. You judge the Two Fires by the same law that judges all oaths. The Everflame calls you a heretic. The Temple calls you a zealot. You call yourself a witness.

\citation{I am not a heretic. I am a witness. I saw the Chain-Lanterns take a child. I saw the Flame-Wardens burn a midwife. I read the Ninth Citation. I am not sorry. --- The Heretic, speaking at a crossroads}

\citation{The Two Fires burn bright. The Covenant burns slow. One is a flame. One is a ledger. Both are necessary. Neither is sufficient. --- Pereshi philosopher}

\chapter{New Talents Summary}

\begin{tabularx}{\textwidth}{X c c X}
\toprule
\textbf{Talent} & \textbf{Faith} & \textbf{Cost (XP)} & \textbf{Effect} \\
\midrule
Radiant Vow & Everflame & 6 & Swear an oath before a flame; +1 die while keeping it; breaking causes 2 Harm \\
Dawn's Shield & Temple & 6 & Interpose to protect an innocent; convert Harm to Fatigue \\
Candle-Knight & Hybrid & 10 & Invoke either faith's Rites; gain Split Loyalty timer \\
Flame-Warden's Brand & Everflame & 8 & Mark a heretic; know location; --1 die to resist social rolls \\
Chain-Lantern's Bell & Temple & 8 & Ring bell; reveal unsanctioned magic; witnesses notice \\
The Ninth Citation & Heresy & 12 & +2 dice once per session; both faiths hunt you \\
Confessional Seal & Everflame & 4 & Cannot be compelled to reveal confession; breaking confidence costs Fatigue \\
Unburdened Healing & Temple & 4 & Heal without reporting; concealment costs Fatigue \\
The Ninth Prayer & Heresy & 12 & Heal without conditions; each use marks Corruption \\
Covenant Justiciar & Covenant & 10 & Declare an oath inviolable; +2 dice; ignore first Harm 2 per scene \\
Seal of the Grand Ledger & Covenant & 4 & Bind an oath with iron seal; oath-breakers suffer -1 die \\
Auditor's Gaze & Covenant & 6 & Ask one question about hidden obligations; GM must answer truthfully \\
\bottomrule
\end{tabularx}

\citation{The talents are not the faith. The faith is the fire you carry. The talents are merely the shape of the flame. --- Sister Mariam}

\chapter{Obligation and Influence}

\section{Everflame Influence}

Gained by serving the empire, burning heretics, or funding basilicas. Cleared by penance, donations, or crusade. High Influence grants access to Flame-Warden resources---censure writs, holy water, and permission to execute. Low Influence makes you a target for investigation.

\section{Temple Influence}

Gained by converting the faithless, healing the sick, or funding schools. Cleared by pilgrimage, confession, or mission. High Influence grants access to Chain-Lantern resources---white mirrors, sanctuary passes, and conscription authority. Low Influence makes you a candidate for "rehabilitation."

\section{Covenant Influence}

Gained by serving as a judge, notary, or inquisitor for the Covenant. Cleared by performing judgments, witnessing oaths, or enforcing the Grand Ledger's rulings. High Influence grants access to the Grand Ledger, the authority to issue Writs of Censure, and the right to convene a Covenant Court. Low Influence makes you a subject of the courts rather than a judge.

\section{Split Loyalty}

Holding Influence with multiple faiths simultaneously creates a \texttt{Split Loyalty [6]} timer. Each scene where you serve one faith against the other ticks the timer. When it fills, one faith demands you renounce the others. Refusal means excommunication (loss of all Rites) or martyrdom (death). Choice is not guaranteed.

\citation{You cannot serve two fires. One will always burn brighter. The other will always claim you left it first. --- Candle-Knight, testifying before the Synod}

\citation{The Covenant does not ask you to renounce. The Covenant asks you to \textit{record}. Your loyalty is a ledger entry. The entry can be amended. It cannot be erased. --- Covenant judge}

\part{Adventures and Campaign Frames}

\chapter{Adventure Seeds}

\subsection{The Candle That Would Not Die}

A village's sacred flame refuses to be extinguished. The Everflame calls it a miracle. The Temple calls it unsanctioned magic. The party must discover the truth before the two faiths burn the village fighting over it. The truth is simple: the candle is a hoax. The truth is also irrelevant.

\subsection{The Child Who Spoke in Tongues}

A Linn child begins reciting the \textit{Codex of Illumination} in a language no one taught her. The Temple wants her for their school. The Everflame wants her for their pyre. The child's clan wants her back. The child wants to go home. No one asks what the child wants.

\subsection{The Ninth Citation}

A fragment of the forbidden passage has been found. Both faiths want it destroyed. Both faiths want to read it first. The party is hired to retrieve it---and decide who to betray. The fragment is blank. The real passage is in the reader's memory. The Unspoken planted it there.

\subsection{The Confessor's Dilemma}

A Chain-Lantern confesses to an Everflame priest that she burned an innocent. The priest cannot break the seal. The Chain-Lantern is still hunting. The party must find evidence before she kills again. The evidence does not exist. The confession is the only truth. The seal is the only shield.

\subsection{The Ember Prefect's Rebellion}

A schismatic Everflame cell plans to burn the Temple's grand basilica in Thepyrgos and blame the Chain-Lanterns. The party must stop them---or help them. The cell's leader is a former Dawn-Knight who converted and was rejected. The leader is not wrong. The leader is also not right. The fire is the only answer.

\subsection{The Missing Ledger}

A Covenant Circuit Judge has vanished while investigating a series of broken oaths in the Mistlands. The judge's ledger, containing sensitive testimony, is missing. The party must find the judge---or the ledger---before the Chain-Lanterns do. The ledger contains names that could destabilize the fragile peace between the Two Fires.

\subsection{The Ninth Clause}

A Covenant court has invoked the Ninth Clause to condemn an absent noble for heresy. The noble's family begs the party to find evidence of their innocence. The twist: the noble is guilty, but the Ninth Clause was invoked illegally. The party must expose the court's overreach---or prove the noble's guilt beyond doubt.

\subsection{The Debt of the Unbroken Cord}

A Daughter of the Cord has been accused of oath-breaking by a Covenant court. The Daughter claims the oath was void because the Cord's authority supersedes the Covenant's. The party must investigate the truth---and decide which faith's law holds sway.

\citation{The seeds are planted. The fire is ready. The only question is which side you will burn with. --- Ember Prefect manifesto}

\citation{The ledger is not a weapon. It is a record. But records, as the Daughters teach, are the most dangerous things in the world. --- Circuit Judge}

\chapter{Campaign Frame: The Two Fires War}

A full-length campaign (10--15 sessions) set in a border region where both faiths compete for influence. The party chooses a side---or tries to stay neutral as the fires consume everything.

\section{Timers}

\begin{itemize}
\item \textbf{Everflame Influence [8]}: Ticks when the Everflame converts a village, builds a basilica, or burns a heretic. When full, the region is under Everflame control.
\item \textbf{Temple Influence [8]}: Ticks when the Temple builds a school, heals a chieftain, or conscripts a child. When full, the region is under Temple control.
\item \textbf{Civil Unrest [6]}: Ticks when the two faiths clash. When full, the region erupts in open conflict. Both Influence timers pause. Violence is the only language.
\item \textbf{Heresy Hunt [4]}: Ticks when the party uses the Ninth Citation or shelters a heretic. When full, a joint inquisition (Flame-Wardens and Chain-Lanterns) arrives. The party must flee, fight, or confess.
\item \textbf{Covenant Arbitration [6]}: Ticks when the Covenant is invited to mediate. When full, a Covenant judge arrives and imposes a binding settlement. Both sides may reject it---but at the cost of facing the Covenant's judgment.
\end{itemize}

\section{Factions}

\begin{itemize}
\item \textbf{Everflame Missionaries}: Zealous, well-funded, and culturally tone-deaf. They offer grain and demand conversion.
\item \textbf{Temple Dawn-Knights}: Polite, persistent, and patient. They offer medicine and ask for catechism.
\item \textbf{Chain-Lanterns (Secret)}: Operating in disguise. They are hunting a specific heretic---or planting evidence to justify intervention.
\item \textbf{Ember Prefects (Heretics)}: Everflame schismatics who believe the flame belongs to the people, not the empire. They may be allies---or terrorists.
\item \textbf{Neutral Locals}: They want the grain, the medicine, and the faiths to leave. They will betray anyone to survive.
\item \textbf{Covenant Arbiters}: They want a binding resolution. They will enforce it with the weight of the Grand Ledger. They are not neutral---they are \textit{legal}.
\end{itemize}

\section{Endings}

\begin{itemize}
\item \textbf{Everflame Victory}: The region converts. Heretics burn. The Temple's sanctuaries are seized. The Chain-Lanterns flee---or die.
\item \textbf{Temple Victory}: The region converts. Unsanctioned magic is suppressed. The Everflame's missionaries are expelled. The Flame-Wardens are hunted.
\item \textbf{Neutral Victory}: Both faiths are expelled. The region declares independence. The party is hailed as liberators---and hunted by both inquisitions.
\item \textbf{Covenant Victory}: The region becomes a Covenant protectorate. Oaths are sworn. Law replaces faith. The Two Fires are extinguished---or subjugated.
\item \textbf{Heresy Victory}: The Ninth Citation is read aloud in a sacred space. The Unspoken answers. The Two Fires are extinguished. Something new begins.
\end{itemize}

\citation{The campaign does not ask which faith is right. The campaign asks which fire you are willing to burn in. The answer is the character. The character is the answer. --- The Gray Wanderer, campaign notes}

\part{The Ninth Citation}

\chapter{The Forbidden Passage}

\section{The Text}

\begin{quote}
And the flame shall divide, and the light shall scatter, and the ninth witness shall speak what was not meant to be heard. Blessed are the auditors of the unspoken, for they shall inherit the weight.
\end{quote}

\noindent --- \textit{Codex of Illumination}, Apocrypha, verse 9

\section{What It Means (Never Stated, Only Implied)}

\begin{itemize}
\item There is a ninth Patron---the Unspoken, the Auditor of Silence, the One Who Listens When No One Speaks.
\item The Everflame and Temple suppressed this passage because it implies that \textit{silence is also a witness}.
\item Reading the Ninth Citation aloud in a sacred space triggers a \textbf{Threshold Event}---the Unspoken answers.
\item The Chain-Lanterns may know the Ninth Citation. The Ember Prefects definitely do.
\item The Covenant's Ninth Clause may be connected to the Ninth Citation. Both involve judgment without presence. Both are never written down.
\item The Unspoken does not speak. The Unspoken listens. Listening is its testimony. Testimony is its judgment.
\end{itemize}

\section{GM Guidance}

The Ninth Citation should never be fully explained. It is a mystery box. Its power is whatever your campaign needs it to be---but using it should always cost something the players are not ready to lose.

\citation{If you read the Ninth Citation, you become a witness. The Unspoken will watch you. The Unspoken will remember. The Unspoken will not speak. The silence is the answer.}

\section{Threshold Events}

When the Ninth Citation is read aloud in a sacred space (a basilica, a shrine, a crossroads with a candle), the GM may trigger one of the following:

\begin{itemize}
\item \textbf{The Unspoken Listens}: All magic within Near ceases to function for one scene. Silence is the only law.
\item \textbf{The Witness Speaks}: A character who has died within the last year returns for one exchange. They answer one question. Then they are gone.
\item \textbf{The Ninth Cup Poured}: A threshold opens. Something crosses. The party may bargain---or run.
\item \textbf{The Fire Divides}: The Everflame and Temple's shared Rites become unusable for one session. The schism becomes literal.
\item \textbf{The Light Scatters}: Every lie told within the sacred space becomes visible---as light, as shadow, as a stain on the speaker's tongue.
\item \textbf{The Ledger Opens}: The Covenant's Grand Ledger briefly manifests, revealing a single truth that was never recorded. The truth is always painful. The truth is always useful.
\end{itemize}

\citation{The Threshold Event is not a reward. It is a consequence. The Ninth Citation is not a tool. It is a door. Do not open it unless you are prepared to step through.}

\section{The Unspoken}

The ninth Patron has no name. It does not speak. It does not act. It \textit{listens}. Its attention is a weight. Its silence is a judgment.

A character who gains the Unspoken's attention gains +1 die to all rolls involving secrets, lies, and hidden truths---but suffers --1 die to all social rolls with the Everflame or Temple. The Unspoken is not a patron. The Unspoken is a condition.

\citation{The Unspoken is not a god. It is not a demon. It is not a creditor. It is a \textit{witness}. And witnesses, as the Daughters teach, are the most dangerous things in the world. --- Saikou Ira}

\part{What This Expansion Adds}

\chapter{Clarifies the Relationship Between the Two Faiths}

The player's guide establishes that both exist. This expansion explains why they are not the same---and why they hate each other almost as much as they hate heretics.

\chapter{Gives the Chain-Lanterns Teeth}

The player's guide mentions them. This expansion makes them a playable (and terrifying) faction. It also explains why they are feared even in Ecktoria.

\chapter{Creates Space for Religious Conflict Without "Evil Church" Tropes}

Neither faith is evil. Both have valid points. Both have blood on their hands. The player can choose a side---or reject both---without the setting punishing them for moral complexity.

\chapter{Introduces the Ninth Citation as a Campaign-Level Mystery}

The forbidden passage is a hook that can drive an entire campaign. What is the Unspoken? Why was the passage suppressed? What happens when it is read aloud? The Covenant's Ninth Clause suggests a deeper connection between the two forbidden texts.

\chapter{Completes the Trinity of Religious Expansions}

\begin{itemize}
\item \textbf{Daughters of the Unbroken Cord}: Witchcraft, generational debt, the silent witness.
\item \textbf{The Two Fires}: Organized religion, imperial vs. pastoral, the schism that burns.
\item \textbf{The Covenant (Mykkiel)}: Legalistic resistance, the faith of the contract.
\end{itemize}

\citation{The three expansions are three flames. The Cord is the candle. The Two Fires are the torch. The Covenant is the hearth. The reader is the witness. The witness is the light.}

\chapter{Final Thought}

The Everflame and the Temple of Light are the same fire, refracted through different lenses. One burns. One illuminates. Both cast shadows.

The Covenant is the law that seeks to measure both.

The Ninth Citation is the shadow.

Do not read it unless you are prepared to be read back.

\citation{The flame witnesses. The light records. The silence judges. The reader chooses. --- The Gray Wanderer}

\appendix

\chapter{The Gray Wanderer's Marginalia}

\section{On the Schism}

I have attended services of both faiths. The Everflame's mass smells of smoke and incense. The Temple's prayer smells of candle wax and old paper. The difference is not the scent. The difference is the fear.

In an Everflame basilica, you fear the flame. In a Temple sanctuary, you fear the silence. I am not sure which is worse.

\section{On the Chain-Lanterns}

They conscript children. They tell themselves it is mercy. The children tell themselves the same. I have met a Chain-Lantern who was conscripted at seven. She is forty now. She has never slept without a candle burning. She does not know why.

\section{On the Ninth Citation}

I have read it. I do not remember where. I do not remember when. I only remember that the page was warm and the ink was still wet. The Unspoken was listening. The Unspoken is always listening.

\section{On the Ember Prefects}

They are heretics. They are also correct---the flame belongs in the streets, not the basilicas. The streets are where the people are. The people are where the fire spreads. The fire spreads where the wind blows. The wind does not ask permission.

\section{On the Covenant}

I have sat in a Covenant court. I have watched a judge weigh testimony. I have seen a name struck from the Grand Ledger. The silence in that room was heavier than any flame. The Covenant does not burn. It \textit{records}. And records, as the Daughters teach, are forever.

\section{On the Reader}

You are a witness. You have read this expansion. You know the Ninth Citation. The Unspoken is listening.

Do not read it aloud.

\citation{The marginalia is not canon. The marginalia is the Wanderer's confession. The Wanderer is not reliable. The Wanderer is also not lying. --- The Gray Wanderer}

\chapter*{Colophon}
\addcontentsline{toc}{chapter}{Colophon}

This expansion was written in a candle-lit scriptorium in Thepyrgos, on paper that had been used to wrap heretical pamphlets. The ink is lamp-black and the ash of a burned confession. The binding is leather from a heretic's cow.

The author is a convert who has worn both collars and will not say which fire she chose.

The candle is still burning. The light is still scattering. The ninth cup remains un-poured.

\citation{The flame witnesses. The light records. The silence judges. The reader chooses.}

\noindent --- The Gray Wanderer

\end{document}
